\documentclass{article}
\usepackage[margin=1in]{geometry}
\usepackage{amsmath, xfrac, actuarialsymbol, mathtools, array}
\usepackage{tabto, xcolor,booktabs,colorPalette}

\DeclarePairedDelimiter\sbrkt{[}{]}
\makeatletter
\let\oldsbrkt\sbrkt
\def\sbrkt{\@ifstar{\oldsbrkt}{\oldsbrkt*}}
\makeatother
\newcommand*{\parens}[2][]{{\left(\vphantom{#1]}#2\right)}}

\definecolor{lander_blue}{RGB}{0,53,153}

\newcount\examUnit
\examUnit=-1 %TODO (based on chapter number)
\begin{document}
  \providecommand{\formulaSheetTitle}[1]{{\Large \textbf{#1}}}
  \newcommand{\quadForm}{
    Quadratic formula: & $\mathllap{x}=\frac{-b\pm\sqrt{b^2-4ac}}{2a}$\\
    Vertex of a parabola: & $\parens{h,k}$ where $h=-\frac{b}{2a}$,\ \  $k=f\parens{h}=c-\frac{b^2}{4a}$\\
  }
  \newcommand{\expoRules}{
    Fractional exponents: & $\mathllap{x^{a/b}}=\sqrt[b]{x^a}$\\
    Subtracting exponents: & $\mathllap{x^{a-b}}=\frac{x^a}{x^b}$\\
  }
  \newcommand{\unitOne}{%
    Average Rate of Change: & $\frac{f(b)-f(a)}{b-a}$\\
    Slope of a Line: & $\mathllap{m}=\frac{y_2-y_1}{x_2-x_1}$\\
    Point-Slope form: & $\mathllap{y-y_1}=m(x-x_1)$\\
    Slope-Intercept form: & $\mathllap{y}=mx+b$\\
  }
  \newcommand{\unitTwo}{%
    General form: & $ax^2+bx+c=0$\\
    Standard form: & $a\parens{x-h}^2+k=0$\\
  }
  \newcommand{\unitThree}{%
    Pythagorean Identity: & $\cos^2(x)+\sin^2(x)=1$\\
    Inverse Trig& $\sin(\theta)=x,\ -\frac{\pi}{2}\leq \theta\leq \frac{\pi}{2}\quad \Longleftrightarrow\quad \theta=\sin^{-1}(x),\ -1\leq x\leq 1$\\
    & $\cos(\theta)=x,\ \phantom{-\frac{\pi}{2}}\mathllap{0}\leq \theta\leq \mathrlap{\pi}\phantom{\frac{\pi}{2}}\quad \Longleftrightarrow\quad \theta=\cos^{-1}(x),\ -1\leq x\leq 1$\\
    & $\tan(\theta)=x,\ -\frac{\pi}{2}< \theta< \frac{\pi}{2}\quad \Longleftrightarrow\quad \theta=\tan^{-1}(x),\ -\infty< x< \infty$\\
  }
  \newcommand{\unitFour}{
    Sum \& Diff: & $\sin(\alpha\pm\beta)=\sin(\alpha)\cos(\beta)\pm \cos(\alpha)\sin(\beta)$\\
    & $\cos(\alpha\pm\beta)=\cos(\alpha)\cos(\beta)\mp \sin(\alpha)\sin(\beta)$\\
    & $\tan(\alpha\pm\beta)=\frac{\tan(\alpha)\pm\tan(\beta)}{1\mp\tan(\alpha)\tan(\beta)}$\\[5pt]
    %
    Double-Angle: & $\sin(2\theta)=2\sin(\theta)\cos(\theta)$ \\
    & $\cos(2\theta)=\cos^2(\theta)-\sin^2(\theta)=1-2\sin^2(\theta)=2\cos^2(\theta)-1$\\
    & $\tan(2\theta)=\frac{2\tan(\theta)}{1-\tan^2(\theta)}$\\[5pt]
    %
    Half-Angle: & $\sin\parens{\frac{\theta}{2}}=
    \pm\sqrt{\frac{1-\cos(\theta)}{2}}$ \hspace*{10mm} $\cos\parens{\frac{\theta}{2}}= \pm\sqrt{\frac{1+\cos(\theta)}{2}}$\\
    & $\tan\parens{\frac{\theta}{2}}= \pm\sqrt{\frac{1-\cos(\theta)}{1+\cos(\theta)}}=\frac{\sin(\theta)}{1+\cos(\theta)}=\frac{1-\cos(\theta)}{\sin(\theta)}$\\
  }
  \begin{center}
    \formulaSheetTitle{Math 114 Formula Sheet}

    \renewcommand{\arraystretch}{2}
    \ifnum\examUnit=-1 \renewcommand{\arraystretch}{1.8} \else\vspace*{\baselineskip} \fi%
    \begin{tabular}{@{}>{\bfseries}p{0.4\linewidth}p{0.59\linewidth}l@{}}
      \ifnum\examUnit=1 \unitOne \fi%
      \ifnum\examUnit=2 \unitTwo \quadForm \fi%
      \ifnum\examUnit=3 \unitThree \fi%
      \ifnum\examUnit=-1 %%Final Exam
        \multicolumn{2}{c}{\hrulefill\textbf{ Unit One }\hrulefill}\\
        \expoRules
        \unitOne
      \multicolumn{2}{c}{\hrulefill\textbf{ Unit Two }\hrulefill}\\
        \unitTwo \quadForm
      \multicolumn{2}{c}{\hrulefill\textbf{ Unit Three }\hrulefill}\\
        \unitThree
      \multicolumn{2}{c}{\hrulefill\textbf{ Unit Four }\hrulefill}\\
        \unitFour
      \fi%
    \end{tabular}
  \end{center}

  \pagebreak
\end{document}